\section{Linear objects, extension data, and derived reconstruction}
\label{t:chapter}

The preceding section enlarged a quotient by retaining identifications and
their coherences.  There is a second enlargement with a different purpose:
keep a nonlinear object, and construct a linear category of coefficients
over it.  Modules then emerge from the construction rather than being
imposed on every object at the outset.  We will see this first for ordinary
groups, then for spaces over $BG$, and finally for rings and their
cotangent complexes.

This distinction also clarifies what Jordan--H\"older does not recover.
Two objects can have the same simple factors and different extension
data.  A derived mapping space can retain the missing extension classes
and their automorphisms.  A derived tensor product can retain a relation
that an ordinary tensor product discards.  These are additional structural
theorems with explicit computations, rather than a claim that a change of
language makes the factors determine the original object.

\subsection{Why a full additive replacement cannot work}

\begin{proposition}\label{t:no-full-additive}
There is no fully faithful functor from the category of groups to an
additive category.
\end{proposition}
\begin{proof}
Suppose $F:\mathbf{Grp}\to\mathcal A$ were such a functor.  The trivial
group has just one endomorphism.  Thus $F1$ has one endomorphism, so its
identity equals its zero endomorphism.  It follows that $F1$ is a zero
object: every map into or out of it is zero.  Consequently $F$ carries
zero homomorphisms to zero morphisms, since these factor through $1$.

Let $p:S_3\to C_2$ be the sign quotient.  There are four homomorphisms
$C_2\to S_3$: the trivial one and the three choices of a transposition
as image of the generator.  There are two homomorphisms $C_2\to C_2$.
Postcomposition with $p$ sends only the trivial map to zero, and sends
the three remaining maps to the identity.  In $\mathcal A$, postcomposition
with $Fp$ is a homomorphism of abelian groups
\[
 \operatorname{Hom}_{\mathcal A}(FC_2,FS_3)
 \longrightarrow \operatorname{Hom}_{\mathcal A}(FC_2,FC_2).
\]
Full faithfulness makes its kernel trivial, so it must be injective.  But
it has the four-element and two-element sets just described as its domain
and codomain.  This is a contradiction.
\end{proof}

The obstruction concerns a \emph{full} replacement preserving all maps.
It does not prohibit a linear invariant, nor an encoding with additional
nonlinear structure.  It explains why we retain the original base when
constructing its linear theory.  The homotopy category of a stable
$\infty$-category is additive, so a fully faithful embedding of the nerve
of all groups into a stable $\infty$-category would give the same
contradiction after passage to homotopy categories.

\subsection{Modules already arise as abelian objects over a group}

Fix a group $G$.  An object of $\mathbf{Grp}/G$ is a homomorphism
$p:E\to G$.  An abelian group object in this slice has, in addition,
a zero section $s:G\to E$, fiberwise addition
\[
 \mu:E\times_G E\longrightarrow E,
\]
and fiberwise inverse, satisfying the abelian group axioms.  All these
maps must themselves be group homomorphisms over $G$.  The requirement
relates the old group multiplication to the new fiberwise addition.

\begin{theorem}[Modules over a group as linear objects]
\label{t:beck}
There is an equivalence of ordinary categories
\[
 \operatorname{Ab}(\mathbf{Grp}/G)
 \simeq \mathbb Z[G]\text{-}\mathbf{Mod}.
\]
\end{theorem}
\begin{proof}
Start with an abelian group object $p:E\to G$ and let $M=\ker p$.
The fiberwise addition restricts to a second operation on $M$ with the
same identity as the inherited group multiplication, since $s$ is a
homomorphism and $s(1)=1$.  The fact that $\mu$ is a homomorphism gives
the interchange law
\[
 \mu(aa',bb')=\mu(a,b)\mu(a',b')\qquad(a,a',b,b'\in M).
\]
It follows that
\[
 \mu(a,b)=\mu(a,1)\mu(1,b)=ab,
 \qquad
 \mu(a,b)=\mu(1,b)\mu(a,1)=ba.
\]
Thus $M$ is abelian, and its two operations agree.

The section acts on $M$ by conjugation:
\[
 g\cdot m=s(g)ms(g)^{-1}.
\]
Normality of the kernel keeps this in $M$.  Conjugation is an automorphism
of its abelian group, and $s(gh)=s(g)s(h)$ gives
$(gh)\cdot m=g\cdot(h\cdot m)$.  Extending the action additively gives
a left $\mathbb Z[G]$-module.  The map
\[
 M\rtimes G\longrightarrow E,\qquad (m,g)\longmapsto m s(g)
\]
is an isomorphism: its inverse sends $e$ to
$(e s(p(e))^{-1},p(e))$.  The multiplication formula is
\begin{equation}\label{t:semidirect-law}
 (m,g)(n,h)=(m+g\cdot n,gh).
\end{equation}
Because $\mu$ is a homomorphism, its values on
$(m s(g),n s(g))$ are obtained by multiplying its value on $(m,n)$ with
its value on $(s(g),s(g))$.  The former is $m+n$ and the latter is $s(g)$,
the zero in the fiber.  Thus the fiberwise addition becomes
$((m,g),(n,g))\mapsto(m+n,g)$.

Conversely, let $M$ be a left $\mathbb Z[G]$-module.  Formula
\eqref{t:semidirect-law} is associative because the action is by
automorphisms and respects multiplication in $G$.  Its identity is $(0,1)$
and the inverse of $(m,g)$ is $(-g^{-1}\cdot m,g^{-1})$.
Projection to $G$ has the zero section $g\mapsto(0,g)$.  Fiberwise addition
is the formula above and fiberwise inverse is $(m,g)\mapsto(-m,g)$.
To check addition is a group homomorphism, compare
\[
 (m+g\cdot m')+(n+g\cdot n')
   =(m+n)+g\cdot(m'+n').
\]
This uses commutativity in $M$ and additivity of the action.  Fiberwise
inverse is a homomorphism by the same additivity, and all abelian group
axioms hold fiber by fiber.

A morphism of abelian group objects preserves $p$, $s$, and $\mu$.
Its restriction to $M$ is additive and commutes with conjugation by
$s(g)$, hence is $\mathbb Z[G]$-linear.  Conversely a module map
$u:M\to M'$ gives $(m,g)\mapsto(u(m),g)$, which preserves every structure
map.  These assignments are inverse on objects up to the displayed
isomorphism and inverse on morphisms.  This proves the equivalence.
\end{proof}

No representation has been chosen externally.  The action appeared
because a coefficient object over $G$ must be transported along $G$.
The theorem derives modules as internal linear objects over a nonlinear
base.  It does not identify a simple $\mathbb Z[G]$-module with a simple
quotient group of $G$; these are different questions about the same base.

\subsection{Stabilization retains the base and enriches the coefficients}

A spectrum object can be described by pointed objects $E_0,E_1,\ldots$
and equivalences $E_n\simeq\Omega E_{n+1}$.  The loop object is formed by
a homotopy pullback at the distinguished point.  Stabilization
$\operatorname{Sp}(\mathcal C)$ collects such objects in a suitable
$\infty$-category $\mathcal C$ with finite limits.  We use the standard
theory in \cite[\S\S1.4.2--1.4.3]{HA}, including the standard $t$-structure on
spectra and its heart $\mathbf{Ab}$.  The heart consists of the
Eilenberg--Mac Lane spectra, having only degree-zero homotopy.

The terminology records a degree-by-degree separation. A spectrum is
connective when $\pi_i=0$ for $i<0$, and coconnective when $\pi_i=0$
for $i>0$. The standard $t$-structure supplies compatible truncations
to these ranges. Its heart is their intersection: only $\pi_0$ can
remain, and it is an abelian group. Thus passing to the heart will
recover ordinary linear coefficients from coefficients with higher degrees.

In the slice $\mathcal S_{/BG}$, a pointed object means a family
$E\to BG$ with a section $BG\to E$; that section chooses the basepoint
in each fiber. Looping is performed in each fiber. A spectrum object
is consequently a family with compatible repeated deloopings, so the
base and its transport remain part of the construction.

\begin{theorem}[Stable coefficients on a groupoid]\label{t:spectra-over-bg}
For an ordinary group $G$ there is a natural equivalence
\[
 \operatorname{Sp}(\mathcal S_{/BG})
  \simeq\operatorname{Fun}(BG,\operatorname{Sp}).
\]
Its pointwise standard $t$-structure has heart
$\mathbb Z[G]\text{-}\mathbf{Mod}$.
\end{theorem}
\begin{proof}
Straightening, Theorem~\ref{h:straightening}, first identifies
$\mathcal S_{/BG}$ with $\operatorname{Fun}(BG,\mathcal S)$.
Pointed objects and finite limits in a functor category are computed
pointwise.  Therefore a sequence of pointed functors with equivalences
$E_n\simeq\Omega E_{n+1}$ is exactly a functor from $BG$ to sequences
of pointed spaces satisfying those equivalences.  This identifies the
objects, coherent maps, and their mapping spaces with
$\operatorname{Fun}(BG,\operatorname{Sp})$.

Pointwise truncation gives the stated $t$-structure.  An object in its
heart assigns an Eilenberg--Mac Lane spectrum to the unique object of
$BG$, hence an abelian group $M$.  Its transport maps give a $G$-action
by automorphisms of $M$.  Mapping spaces between heart objects are
discrete, so the coherent conditions reduce exactly to the ordinary
action identities.  Maps are additive maps commuting with the action.
These are left $\mathbb Z[G]$-modules and their linear maps.
\end{proof}

This theorem adds higher coefficient groups and coherent transport while
keeping $BG$ in the construction.  Its restriction to degree zero has
the same module category as Theorem~\ref{t:beck}.  The two descriptions
explain the relationship between ordinary internal linear objects and
stable local coefficient systems.

The ambient homotopical category matters.  We have stabilized
$\mathcal S_{/BG}$, not the ordinary nerve of $\mathbf{Grp}$.
Indeed, in the ordinary pointed category of groups the pullback
$1\times_G1$ is $1$ for every $G$.  Thus its ordinary loop construction
is trivial, and spectrum objects in that unenhanced nerve are trivial.
The homotopy fiber $*\times^h_{BG}*$ instead retains $G$.
Homotopical enhancement must precede this stabilization.

\subsection{A derived connecting map reconstructs an extension}

For this subsection $R$ is an associative unital ring and modules are left
modules.  We use cochain complexes, whose differentials raise degree.
The shift is $(C[1])^j=C^{j+1}$ with negated differential, so a module
$N[1]$ has its only nonzero cohomology in degree $-1$.

We use the standard derived framework: localizing complexes at
quasi-isomorphisms gives a stable $\infty$-category $\mathcal D(R)$;
its heart is $R\text{-}\mathbf{Mod}$; a fiber/cofiber sequence gives the
long exact cohomology sequence; and a bounded projective resolution
computes derived maps and derived tensor products.  These are foundational
inputs, not consequences of composition-factor uniqueness.  References
are \cite[\S1.3]{HA} and \cite[Tags 06XP, 0643]{Stacks}.
Write $[X,Y]$ for $\pi_0\operatorname{Map}_{\mathcal D(R)}(X,Y)$ and set
\[
 \operatorname{Ext}^1_R(L,N)=[L,N[1]].
\]

\begin{theorem}[Extensions are connecting morphisms]
\label{t:extension-classification}
The isomorphism classes of marked short exact sequences
\[
 0\longrightarrow N\xrightarrow{i}E\xrightarrow{p}L\longrightarrow0
\]
are naturally in bijection with $[L,N[1]]$.  Here an isomorphism of marked
extensions induces the identity on both $N$ and $L$.
\end{theorem}
\begin{proof}
The cone of $i:N[0]\to E[0]$ is the complex $N\to E$ in degrees
$-1,0$.  Its cohomology is $L$ in degree zero and zero elsewhere, so
the quotient map identifies it with $L$ in $\mathcal D(R)$.  The cofiber
sequence therefore supplies a connecting map
\[
 N\longrightarrow E\longrightarrow L\xrightarrow{\delta}N[1].
\]
An isomorphism of marked extensions preserves $\delta$.

Conversely, given a map $\delta:L\to N[1]$, form its fiber $E_\delta$.
The long exact cohomology sequence gives
\[
 0\longrightarrow N\longrightarrow H^0(E_\delta)
 \longrightarrow L\longrightarrow0,
\]
and $H^j(E_\delta)=0$ for $j\neq0$.  The canonical truncation maps
therefore identify $E_\delta$ with this module in the derived category.
A path between two choices of $\delta$ induces an equivalence of their
fiber diagrams, hence an isomorphism of marked extensions.

Finally, the fiber of the connecting map of the first construction is
the original $E$, with its given maps.  The connecting map recovered
from the second construction is the original $\delta$, because both
are the same fiber/cofiber diagram.  This uses the universal properties
of the fiber and the stable rotation of a cofiber sequence; the spaces
of choices of each such limit or colimit are contractible.  Thus the
two constructions are inverse on the stated isomorphism classes.
\end{proof}

\begin{proposition}[The automorphisms also survive]
\label{t:extension-automorphisms}
The automorphism group of a marked extension as above is naturally
$\operatorname{Hom}_R(L,N)$.  It agrees with the fundamental group of
$\operatorname{Map}(L,N[1])$ at its class $\delta$.
\end{proposition}
\begin{proof}
For $u:L\to N$, the map $1_E+iup$ preserves both markings and has
inverse $1_E-iup$, since $pi=0$.  Multiplication of these automorphisms
corresponds to addition of $u$.
Conversely, if $f$ preserves the markings, $p(f-1)=0$, so $f-1$ factors
uniquely as $iv$ with $v:E\to N$.  The equality $(f-1)i=0$ and
monicity of $i$ give $vi=0$.  The quotient property of $p$ makes $v=up$
for a unique $u:L\to N$.

In a stable category, looping the mapping space at zero identifies
$\pi_1\operatorname{Map}(L,N[1])$ with $[L,N]$.  Translation in this
grouplike mapping space gives the same group at $\delta$.
The fully faithful inclusion of the heart identifies $[L,N]$ with
$\operatorname{Hom}_R(L,N)$, proving the final assertion.
\end{proof}

The bijection also respects the Baer addition previously
constructed. Naturality of the connecting map makes pullback along
$L\to L\oplus L$ and pushout along $N\oplus N\to N$ act by
precomposition and postcomposition. Applied to $\delta_1\oplus\delta_2$,
their composite is $\delta_1+\delta_2$. Thus the derived notation
$\operatorname{Ext}^1$ identifies the same abelian group as the earlier
extension construction, not only the same underlying set.

The factors $N,L$ are the endpoints, and $\delta$ is the missing
attachment.  Its mapping space also retains the automorphisms of that
attachment.  This is a precise instance of gaining reconstruction data
by retaining maps rather than just a list of objects.

\begin{example}[The same two factors, different attachments]
\label{t:dual-number-extension}
Let $R=k[t]$ for a field $k$, and let $k=R/(t)$ as an $R$-module.
The free resolution
\[
 0\longrightarrow R\xrightarrow{\,t\,}R\longrightarrow k
 \longrightarrow0
\]
is exact: multiplication by $t$ is injective in the polynomial ring,
and its image is the kernel of evaluation at $0$.  Applying
$\operatorname{Hom}_R(-,k)$ makes the differential zero, hence
$\operatorname{Ext}^1_R(k,k)\cong k$.

Every marked extension can also be described without derived notation.
Let $e$ be the specified image of $1$ in its kernel and choose a lift $f$
of $1$ in its quotient.  Then $e,f$ span the underlying vector space:
subtract a multiple of $f$ from any element and the result is in $ke$.
They are independent by applying the quotient map first, then using
injectivity on the kernel.  Necessarily
\[
 te=0,\qquad tf=\lambda e
\]
for a unique $\lambda\in k$.  Replacing $f$ by $f+ce$ does not change
$\lambda$.  Conversely these formulas define an $R$-module extension
for every $\lambda$.  A map preserving the markings preserves this
scalar, and the formulas give an isomorphism whenever the scalars agree.
This is the concrete parameter detected by the derived calculation.

At $\lambda=0$ the module is $k\oplus k$.  At $\lambda\neq0$ it is
isomorphic, after forgetting the markings, to $R/(t^2)$: send $1$ to
$f$ and $t$ to $\lambda e$.  The map is surjective, has kernel $(t^2)$,
and the displayed independent vectors prove injectivity of the induced
map.  Both modules have two factors isomorphic to $k$, but their
$t$-operators are different.  The marked classification distinguishes
all values of $\lambda$, while the middle-object classification merges
all nonzero values.  These classifications must not be conflated.
\end{example}

\subsection{A derived tensor product retains an equation}

For commutative $R$, the derived tensor product is denoted
$M\otimes_R^{\mathbf L}N$.  A projective resolution replaces one input
before tensoring.  This repairs the failure of ordinary tensoring to
preserve exactness by retaining the resulting homology.

\begin{proposition}\label{t:derived-intersection}
For $R=k[t]$ and $k=R/(t)$,
\[
 H^j(k\otimes_R^{\mathbf L}k)=
 \begin{cases}k,&j=0,-1,\\0,&\text{otherwise.}\end{cases}
\]
The ordinary tensor product is only $k$.
\end{proposition}
\begin{proof}
Place the two free terms of the resolution in
Example~\ref{t:dual-number-extension} in degrees $-1,0$.
Tensoring with $k$ gives the complex
\[
 \bigl[\,k\xrightarrow{\,0\,}k\,\bigr]
\]
in those degrees, because $t$ acts as zero on $k$.
Its cohomology is exactly the displayed answer.
Ordinary tensoring gives $R/(t)\otimes_R R/(t)\cong R/(t)$: the map
$a\otimes b\mapsto ab$ has inverse $c\mapsto c\otimes1$.
\end{proof}

In affine geometry this is the derived self-intersection of the origin
inside the affine line.  The degree $-1$ copy of $k$ retains the overlap
of the two imposed copies of the equation $t=0$.  The statement about
cohomology was proved directly by the resolution; the geometric name
does not serve as its proof.  The calculation is already beyond what a
list of simple factors or an ordinary quotient alone records.

\subsection{Cotangent objects: the same method for equations and deformations}

There is a corresponding linear theory for a commutative algebra.
For an ordinary $A$-module $M$, the split square-zero algebra $A\oplus M$
has multiplication
\[
 (a,m)(b,n)=(ab,an+bm).
\]
Its ideal $0\oplus M$ has square zero.  A section of
$A\oplus M\to A$ has the form $a\mapsto(a,d(a))$.  Comparing products
shows that it is a ring map exactly when
$d(ab)=a\,d(b)+b\,d(a)$ and $d$ is additive with $d(1)=0$.
Over a base ring $R$, being an $R$-algebra map adds $d(r)=0$.
Thus derivations arise as maps into linear objects over the algebra.

For the higher statement, the notation retains these markings:
\[
 \operatorname{Der}_A(B,M)
 :=\operatorname{Map}_{(\operatorname{CAlg}_{A/})_{/B}}
       (B,B\oplus M).
\]
The source is the identity map of the $A$-algebra $B$ and the target is
its split square-zero extension, with its projection to $B$.  Thus this
is the space of sections over $B$ in $A$-algebras, including the specified
coherences with the base and the projection.

\begin{theorem}[Framework theorem, quoted: tangent modules and cotangents]
\label{t:cotangent-foundation}
For an $E_\infty$-ring $A$,
$\operatorname{Sp}(\operatorname{CAlg}_{/A})\simeq\operatorname{Mod}_A$.
For a map $A\to B$, derived $A$-derivations of $B$ into a $B$-module $M$
are represented by a cotangent module $L_{B/A}$:
\[
 \operatorname{Der}_A(B,M)
 \simeq\operatorname{Map}_{\operatorname{Mod}_B}(L_{B/A},M).
\]
For $A\to B\to C$ there is a cofiber sequence
\begin{equation}\label{t:transitivity}
 C\otimes_B L_{B/A}\longrightarrow L_{C/A}
 \longrightarrow L_{C/B}.
\end{equation}
\end{theorem}
These are the tangent equivalence, representing property, and transitivity
formalism of \cite[\S7.3; Corollaries 7.3.4.14 and 7.3.3.6]{HA}.
Tensor products here are tensor products of module spectra, hence derived.
In ordinary commutative algebra the corresponding cotangent complexes
and transitivity triangle are constructed using simplicial resolutions
\cite[Tag 08QX]{Stacks}.

The category $\operatorname{CAlg}_{/A}$ contains nonlinear algebras mapping
to $A$.  Its stabilization produces the linear coefficient category at
$A$, just as spaces over $BG$ produced spectrum-valued local coefficients.
The representing property says exactly which measurements the cotangent
module encodes: all derived derivations into any coefficient module.

\begin{proposition}[An explicit obstruction to extending a derivation]
\label{t:derivation-obstruction}
Let $A\to B\to C$ be maps of $E_\infty$-rings and $M$ a $C$-module.
A homotopy class of $A$-derivations $B\to M$ extends to an $A$-derivation
$C\to M$ if and only if its image under the connecting homomorphism
\[
 [C\otimes_B L_{B/A},M]
   \longrightarrow [L_{C/B},M[1]]
\]
vanishes.  The brackets denote morphisms in the homotopy category of
$C$-modules.
\end{proposition}
\begin{proof}
Apply the mapping-spectrum functor into $M$ to
\eqref{t:transitivity}.  Mapping out of a cofiber turns it into a fiber,
giving the exact sequence
\[
 [L_{C/A},M]\longrightarrow[C\otimes_B L_{B/A},M]
   \longrightarrow[L_{C/B},M[1]].
\]
The first arrow is restriction of derivations: naturality of the
representing property identifies the cotangent map with restriction,
and extension of scalars identifies the middle term with
$[L_{B/A},M]$ as a $B$-module mapping group.
Exactness says precisely that a class lies in the image of restriction
if and only if its connecting class is zero.
\end{proof}

This is a theorem about existence of an extension of a homotopy class.
It does not assert uniqueness.  The full mapping fiber retains the
choices of extension, their paths, and their higher compatibilities.
The transitivity triangle organizes all three algebras in one statement;
its connecting arrow identifies the obstruction rather than merely
asserting that some obstruction ought to exist.

For a final explicit computation, use ordinary commutative algebras over
a field $k$ of characteristic zero.  This choice also permits comparison
with commutative differential graded algebras and $E_\infty$-$Hk$-algebras
\cite[Proposition 7.1.4.11]{HA}; no equivalence between those models in
arbitrary positive characteristic is being assumed. The calculation below
is in ordinary derived cotangent theory. Under the rational comparison,
its cotangent object corresponds to the spectral one because both
corepresent the same derived-derivation functor, with the algebra and
coefficient module retained.

\begin{theorem}[Framework theorem, quoted: a hypersurface cotangent]
\label{t:hypersurface-input}
If $P=k[x_1,\ldots,x_n]$ and $A=P/(f)$ with $f$ a nonzero divisor,
the ordinary derived cotangent complex $L_{A/k}$ is represented in
degrees $-1,0$ by
\[
 (f)/(f^2)\longrightarrow\Omega_{P/k}\otimes_P A,
 \qquad [f]\longmapsto df.
\]
\end{theorem}
This regular-quotient computation follows from the cotangent
transitivity triangle and the regular-sequence theorem
\cite[Tags 08SH, 08QX]{Stacks}.  The ordinary module of differentials
alone is its degree-zero cohomology; the complex retains the relation
before taking that quotient.

\begin{example}[The infinitesimal direction of the dual numbers]
\label{t:cotangent-dual-numbers}
Take $A=k[t]/(t^2)$.  The polynomial $t^2$ is a nonzero divisor in
$k[t]$, so Theorem~\ref{t:hypersurface-input} applies.  Both terms are
free rank-one $A$-modules, and the complex is
\[
 L_{A/k}\simeq
 \bigl[\,A e\xrightarrow{\ e\mapsto 2t\,dt\ }A\,dt\,\bigr].
\]
The characteristic assumption makes $2$ invertible.  Multiplication by
$2t$ has image $kt$ and kernel $kt$ in $A$.  Consequently
\[
 H^0(L_{A/k})\cong k\,dt,
 \qquad H^{-1}(L_{A/k})\cong k\,(te).
\]
These assertions can be checked on $a+bt$: multiplication by $2t$
gives $2at$.  The negative-degree term therefore survives; replacing
the complex by its ordinary differentials loses it.

Let $k$ be an $A$-module through $t\mapsto0$.  Since the displayed
cotangent complex is a bounded complex of free modules, its derived
maps into $k$ are computed by the ordinary Hom complex.  That complex
has a copy of $k$ in degrees $0,1$ and zero differential.  Thus
\[
 [L_{A/k},k]\cong k,
 \qquad [L_{A/k},k[1]]\cong k.
\]
The first is the group of derivations into this coefficient module;
the second is the next group available for obstruction calculations
such as Proposition~\ref{t:derivation-obstruction}.  A nonzero possible
obstruction group does not by itself assert that a specified lifting
problem has a nonzero obstruction: that requires computing its
connecting class.
\end{example}

The progression is now visible.  Exact quotients and refinement identify
which simple pieces are invariant.  Coherent quotients retain actions and
transport.  Linear objects over a base produce modules.  Derived maps
retain attachments and automorphisms; derived tensor products and
cotangent complexes retain relations and obstruction classes.  Each
construction answers a new question about information the preceding
one did not retain.  The ordinary Jordan--H\"older theorem remains one
proved consequence of the exact framework, while these higher and
derived results enlarge the kinds of reconstruction that can be asked
and answered.
